\begin{abcliste}
\abc Beim Fallen wird die potentielle Energie (Lageenergie) der Kugel in kinetische Energie umgewandelt. Beim Aufprall werden Kugel und Boden kurz elastisch verformt: Die kinetische Energie wird zu Spannenergie, die beim Entspannen wieder in kinetische Energie zurückverwandelt wird. Beim Aufsteigen wird die kinetische Energie wieder zu potentieller Energie. Ein kleiner Teil der Energie wird bei jedem Aufprall in innere Energie (Erwärmung) und Schall umgewandelt; deshalb erreicht die Kugel nicht ganz die Ausgangshöhe.
\abc Beim Anlauf wird chemische Energie der Muskeln in kinetische Energie umgewandelt. Auf dem Minitrampolin wird diese in Spannenergie des Trampolins umgewandelt, die den Turner wieder nach oben beschleunigt (kinetische Energie, ergänzt durch Muskelarbeit beim Absprung). Beim Steigen wird kinetische Energie zu potentieller Energie; ein Teil bleibt als kinetische Energie der Vorwärtsbewegung und der Drehung erhalten. Beim Fallen auf die Matte wird die potentielle Energie wieder zu kinetischer Energie, die beim Landen durch die Verformung der Matte vollständig in innere Energie (Wärme) und Schall umgewandelt wird.
\end{abcliste}
